\chapter{Pipeline sperimentale e metodologia}

\section{Dataset e task}

Per realizzare la tesi sono stati utilizzati due dataset:

\begin{itemize}
	\item \emph{Oxford-IIIT Pet}~\cite{parkhi12a}, composto da immagini di cani e gatti suddivise in 37 classi.
	      Il dataset comprende un file di annotazione che associa a ogni immagine una tupla
	      (\texttt{image\_path}, \texttt{micro\_label}, \texttt{macro\_label}), ove:

	      \begin{itemize}
		      \item \texttt{micro\_label} identifica la razza tra le 37 disponibili;
		      \item \texttt{macro\_label} identifica la specie, ovvero cane o gatto.
	      \end{itemize}

	      Il dataset è diviso in tre sottoinsiemi: training (70\%), validazione (15\%) e test (15\%).
	      La suddivisione è effettuata in modo da preservare l'equa distribuzione delle 37 classi.
	\item \emph{CIFAR-100}~\cite{krizhevsky2009learning}, composto da 100 classi raggruppate in 20 superclassi.
	      Ogni superclasse comprende 5 classi, per un totale di 100 \textit{fine label}. Per contenere il costo
	      computazionale degli esperimenti, è stato estratto inizialmente un sottoinsieme di 50 campioni per ciascuna classe.
	      Successivamente tale numero è stato aumentato a 200 campioni per classe, poiché la quantità iniziale produceva un risultato
	      caratterizzato da variabilità e rumore.

	      Il campionamento è stato effettuato prima della suddivisione stratificata del dataset, secondo la proporzione 70\%/15\%/15\% tra training, validazione e test.
\end{itemize}

\section{Architettura multi-task}
\label{chapter:3.2}

Delle famiglie architetturali presentate nella Sezione~\ref{chapter:2.2}, sono state selezionate quattro varianti come backbone per gli esperimenti: \textbf{ResNet18}, \textbf{ResNet50}, \textbf{DenseNet121} ed \textbf{EfficientNet-B0}.

L'integrazione dell'apprendimento multi-task su tali backbone segue uno schema standard, noto come \emph{hard parameter sharing}: un unico backbone convoluzionale condiviso è utilizzato da due head di classificazione indipendenti.

Il backbone elabora l'immagine fino al GAP\footnote{GAP: \emph{Global Average Pooling}. Operazione applicata al termine del backbone che riduce ciascuna mappa di feature a un singolo valore medio, trasformando così il volume di feature tridimensionale in un vettore monodimensionale utilizzabile dalle head di classificazione.}, producendo un vettore di embedding la cui dimensione dipende dall'architettura:

\begin{itemize}
	\item 512 per ResNet18;
	\item 2048 per ResNet50;
	\item 1024 per DenseNet121;
	\item 1280 per EfficientNet-B0.
\end{itemize}

Questo vettore è utilizzato in parallelo da entrambe le head.

Le due head sono costituite da un singolo layer lineare (\texttt{Linear(feature\_dim, num\_classes)}). Le proiezioni delle due head sono, rispettivamente, su 37 classi (razza) e 2 classi (specie) per Oxford-IIIT Pet; su 100 classi (\emph{fine}) e 20 classi (\emph{coarse}) per CIFAR-100.

La funzione di loss è una combinazione pesata delle due cross-entropy, tramite i coefficienti $\lambda_{\text{micro}}$ e $\lambda_{\text{macro}}$:

\begin{equation}
	\mathcal{L}_{\text{tot}} = \lambda_{\text{micro}} \cdot \mathcal{L}_{\text{micro}} + \lambda_{\text{macro}} \cdot \mathcal{L}_{\text{macro}},
\end{equation}

con $\lambda_{\text{micro}} = 1.0$ e $\lambda_{\text{macro}} = 0.5$: il task \emph{fine} (micro) pesa dunque il doppio rispetto al task \emph{coarse} (macro) nella loss totale.

\section{Procedura di training dei modelli}
\label{sec:procedura-training}

\subsection{Regime di addestramento}

Tutti i backbone sono inizializzati con i pesi pre-addestrati su ImageNet, secondo l'approccio del \emph{transfer learning}~\cite{yosinski2014transferablefeaturesdeepneural}. 
Come mostrato dagli stessi autori, la specializzazione dei layer più profondi al task originario può limitare la trasferibilità delle feature se queste non vengono adattate al nuovo dominio; 
per questo motivo l'addestramento è condotto in modalità \emph{full fine-tuning}: tutti i parametri del backbone rimangono aggiornabili, 
insieme a quelli delle due head, anziché congelare il backbone e addestrare solo le head (\emph{feature extraction}).

Il backbone adatta così i propri filtri alle caratteristiche peculiari del dataset, invece di limitarsi alle feature generiche di ImageNet.

\subsection{Ottimizzazione}

I pesi sono aggiornati tramite l'ottimizzatore \emph{Adam}, con learning rate fissato a $0.001$, minimizzando la loss combinata $\mathcal{L}_{\text{tot}}$ definita nella Sezione~\ref{chapter:3.2}. L'addestramento è condotto per un massimo di 50 epoche, con dimensione del batch pari a 32.

\subsection{Trasformazioni sui dati}

Tutte le immagini sono ridimensionate a $224 \times 224$ pixel e normalizzate secondo media e deviazione standard di ImageNet ($\mu = [0.485, 0.456, 0.406]$, $\sigma = [0.229, 0.224, 0.225]$), 
ovvero gli stessi valori usati durante il pre-training su ImageNet.

Nel set di training, sull'immagine ridimensionata è applicato un \emph{random horizontal flip} con probabilità $p = 0.5$, come unica forma di data augmentation. Nei set di validazione e test non è applicata alcuna trasformazione stocastica: le immagini sono unicamente ridimensionate e normalizzate, per garantire una valutazione deterministica e riproducibile.

Tutte le operazioni stocastiche (suddivisione del dataset, degradazione controllata, augmentation) sono governate da un unico seed fisso
che garantisce il determinismo\footnote{L'esperimento si può riprodurre molteplici volte ottenendo lo stesso risultato, grazie al seed fissato come parametro.} dell'intera pipeline.

\subsection{Early stopping}

Il training è monitorato a ogni epoca, valutando la loss totale $\mathcal{L}_{\text{tot}}$ (calcolata con gli stessi coefficienti $\lambda_{\text{micro}}$ e $\lambda_{\text{macro}}$ definiti in Sezione~\ref{chapter:3.2}) sul set di validazione. 
È adottato un criterio di \emph{early stopping} con pazienza pari a 6 epoche: se la loss di validazione non migliora per 6 epoche 
consecutive, l'addestramento viene interrotto anticipatamente, ripristinando i pesi del modello relativi all'epoca in cui si è 
registrata la loss di validazione minima, corrispondente al miglior modello.

\section{Degradazione controllata del training set}
\label{sec:degradazione}

Per studiare la robustezza delle architetture a diverse forme di perdita di informazione, 
è stata implementata una procedura di degradazione controllata del training set, articolata in due meccanismi indipendenti: \emph{macro-drop} e \emph{micro-drop}. 
Entrambi sono applicati esclusivamente al training set, lasciando i set di validazione e test sempre integri. 
Le percentuali di degradazione applicate vanno dallo 0\% al 40\%, con incrementi del 5\%, e l'intera procedura è determinata da un seed.

\subsection*{Macro-drop}

Il \emph{macro-drop} simula una condizione di scarsità di dati, riducendo il numero di campioni disponibili senza mai eliminare interamente le classi. 
La rimozione è stratificata per classe: per ogni classe viene rimossa la medesima quantità proporzionale di campioni; eventuali resti derivanti da arrotondamento vengono distribuiti in modo che alcune classi perdano 
un campione in più o in meno rispetto alle altre. 

Tutte le classi restano quindi rappresentate nel training set, con un numero di esempi ridotto. Poiché la rimozione avviene sull'intero training set, il \emph{macro-drop} incide contemporaneamente su entrambe le head del modello.

\subsection*{Micro-drop}

Il \emph{micro-drop} simula la scomparsa di intere classi durante il training. A ogni percentuale di degradazione, un sottoinsieme di classi (razze per Oxford-IIIT Pet, classi \emph{fine} per CIFAR-100) viene selezionato ed eliminato dal training set, il quale corrisponde al 70\% del dataset totale.

Per Oxford-IIIT Pet, la selezione è ristretta alle sole razze canine (circa 3484 immagini nel training set, contro le 1659 delle razze feline), mentre per CIFAR-100 non è necessaria alcuna restrizione: le classi da eliminare vengono scelte casualmente, tramite seed, tra tutte le 100 classi \textit{fine}. La scelta di filtrare Oxford-IIIT Pet nasce dallo sbilanciamento del dataset tra le due specie (25 razze canine contro 12 feline): una selezione casuale su tutte le 37 razze avrebbe potuto, in alcuni casi sfortunati, colpire in prevalenza o esclusivamente le razze feline, già numericamente sottorappresentate, alterando in modo incontrollato l'equilibrio tra le due categorie di livello superiore (specie).

Il numero di classi da eliminare è determinato a partire dalla percentuale di drop richiesta, calcolata sul numero di immagini anziché sul numero di classi. Ad esempio, un drop del 20\% su 3484 immagini corrisponde a circa 696 immagini da rimuovere; poiché ogni razza dispone di circa 140 immagini nel training set, questo equivale a circa 4.97 classi. L'algoritmo elimina le classi in ordine progressivo, secondo il seed fisso, rimuovendo per intero ciascuna classe prima di passare alla successiva: nell'esempio, le prime 4 classi vengono eliminate completamente, mentre la quinta viene svuotata solo parzialmente, fino al raggiungimento del numero target di immagini.

Ad esempio, con \emph{micro-drop} al 20\% su Oxford-IIIT Pet, applicato con seed \texttt{777}, sono state selezionate e rimosse le razze \emph{Yorkshire Terrier}, \emph{English Setter}, \emph{German Shorthaired Pointer} e \emph{Boxer}. 


La head relativa al task coarse (macro, ovvero specie o superclasse) continua ad addestrarsi regolarmente sui campioni delle classi superstiti: poiché la selezione delle istanze da escludere opera su singole categorie di livello fine (micro), le restanti classi appartenenti alla medesima superclasse continuano a fornire esempi per l'aggiornamento della head coarse.

\subsection*{Effetto sulla valutazione}

Indipendentemente dal tipo e dall'entità della degradazione applicata in training, la valutazione finale è sempre condotta sull'insieme completo delle classi (37 per Oxford-IIIT Pet, 100 per CIFAR-100). 
Con il \emph{micro-drop}, questo significa che il modello viene valutato anche su classi che non ha mai osservato durante l'addestramento. Tale condizione ha conseguenze dirette sulle metriche di valutazione, 
discusse nella Sezione~\ref{sec:metriche-valutazione}.

\section{Metriche di valutazione}
\label{sec:metriche-valutazione}

I risultati degli esperimenti sono raccolti in tre tipologie di file CSV, con scopi complementari.

\subsection*{Log di training}

Per ogni modello viene generato un unico file di log (ad esempio \texttt{resnet18\_training\_log.csv}) che aggrega lo storico di tutte le run effettuate su quel backbone. 
Ogni riga corrisponde a un'epoca di una specifica run, identificata dal tipo di degradazione, dalla percentuale e dal seed, e riporta le corrispondenti loss di training e di validazione.

\subsection*{Report globale}

Il report globale presenta, per ogni combinazione di tipologia di degradazione, percentuale e seed, tre valori aggregati calcolati sul test set completo (37 classi per Oxford-IIIT Pet, 100 per CIFAR-100):

\begin{itemize}
    \item \textbf{Accuracy}: rapporto tra predizioni corrette e totale delle predizioni, calcolata sul task micro;
    \item \textbf{Macro F1 globale}: media aritmetica dell'F1-score calcolato su tutte le classi del task micro, incluse quelle non presenti nel training della run corrente;
    \item \textbf{Macro F1 note}: media aritmetica dell'F1-score calcolato esclusivamente sulle classi effettivamente presenti nel training della run corrente. 
	L'insieme delle classi note viene ricostruito a partire dal seed della run, ripetendo la stessa procedura di degradazione applicata in fase di addestramento.
\end{itemize}

La distinzione tra le due varianti di F1 è motivata dall'effetto del \emph{micro-drop} sulle metriche: le classi escluse dal training ricevono soltanto gradienti che abbassano il logit della loro unità di output corrispondente e, 
in fase di test, ottengono tipicamente F1 pari a zero (nessuna predizione corretta per quella classe). 

Includendo queste classi nel calcolo della Macro F1 globale, la metrica risulta penalizzata dalle classi escluse, 
indipendentemente dalla reale capacità del modello sulle classi effettivamente apprese.
 
La \emph{Macro F1 note} isola questo effetto, restituendo una misura della performance del modello limitata a ciò che ha potuto effettivamente imparare.

\subsection*{Report per classe}

Il terzo file riporta, per ogni combinazione di tipologia di degradazione, percentuale e seed, una riga per ciascuna delle classi del task micro, con le seguenti informazioni:

\begin{itemize}
    \item l'indicazione se la classe sia stata osservata o meno durante il training della run corrente;
    \item precision, recall, F1-score e support, calcolati sul test set.
\end{itemize}

Questo livello di dettaglio consente di analizzare l'impatto della degradazione classe per classe, distinguendo il comportamento del modello sulle classi note da quello sulle classi mai osservate in training.